\begin{figure}[htbp]
\centering
\begin{tikzpicture}[
    st/.style={base, text width=2.42cm, minimum height=1.55cm},
  ]
  \def\dx{2.92}
  % top row, left to right
  \node[st, n]     (lit)  at (0*\dx,0) {\textbf{Evidence}\\ source papers,\\ state of the art\\ \textit{Ch.~3--4}};
  \node[st, alert] (gap)  at (1*\dx,0) {\textbf{Gap}\\ saturated and\\ open directions\\ \textit{Ch.~5}};
  \node[st, fuse]  (meth) at (2*\dx,0) {\textbf{Method}\\ \method, \methodS,\\ hypotheses\\ \textit{Ch.~6}};
  \node[st, dat]   (data) at (3*\dx,0) {\textbf{Data}\\ core stack,\\ alignment audit\\ \textit{Ch.~7}};
  \node[st, n]     (prot) at (4*\dx,0) {\textbf{Protocol}\\ metrics, baselines,\\ ablations\\ \textit{Ch.~8}};

  % bottom row, right to left
  \node[st, tim]   (phas) at (4*\dx,-2.75) {\textbf{Build}\\ Phases 0--4,\\ notebooks\\ \textit{Ch.~9--10}};
  \node[st, fuse]  (res)  at (3*\dx,-2.75) {\textbf{Results}\\ tables built\\ from run logs\\ \textit{Notebook 15}};
  \node[st, n]     (eth)  at (2*\dx,-2.75) {\textbf{Ethics}\\ privacy review,\\ approvals\\ \textit{Ch.~11}};
  \node[st, alert] (pap)  at (1*\dx,-2.75) {\textbf{Manuscript}\\ journal choice,\\ checklist\\ \textit{Ch.~12}};

  \draw[arr] (lit) -- (gap);
  \draw[arr] (gap) -- (meth);
  \draw[arr] (meth) -- (data);
  \draw[arr] (data) -- (prot);
  \draw[arr] (prot) -- (phas);
  \draw[arr] (phas) -- (res);
  \draw[arr] (res) -- (eth);
  \draw[arr] (eth) -- (pap);

  % feedback loops
  \draw[darr] (res.north west) -- node[lbl, fill=white, pos=0.5] {revise} (meth.south east);
  \draw[darr] (pap.north) -- node[lbl, fill=white, pos=0.5] {re-check} (gap.south);
\end{tikzpicture}

\vspace{2mm}
\caption[The research programme and its feedback loops]{The research
programme. Solid arrows give the order of work and the chapters that
describe each stage. The two dashed arrows are the feedback loops that
matter most. Ablation results revise the method. Prior art is checked again
immediately before submission, because the field publishes new Mamba
detectors every month.}
\label{fig:programme}
\end{figure}
